\chapter{子流形、向量场与可积分布}
\label{ch:05-submanifolds}
\ChapterPlan{周1.3至1.5（本地页码待核对）；梁上册第二版2.2.2，印刷页27--30，PDF物理页43--46；其余为自编补充}{正文与证明初稿；附8题解答及原创图，尚未独立数学审阅}

第4章给出了逐点的一阶工具。本章先用它识别方程的解集，再研究向量场如何产生运动，最后判断一族切平面能否拼成局部曲面。依赖顺序为常秩定理、ODE局部理论、流盒、李括号、Frobenius局部坐标定理；不提前使用单位分解或微分形式的外微分。

\section{浸入、嵌入与正则值}
线性映射的单射、满射分别刻画参数是否冗余与方程是否独立。将这些条件施加于每一点的切映射，便得到浸入和浸没；嵌入还必须控制整体拓扑。
\begin{definition}[浸入、浸没与嵌入]
设$F:M^m\to N^n$光滑。若每个$\dd F_p$单射，称$F$为浸入（immersion）；若每个$\dd F_p$满射，称为浸没（submersion）。若$F$为浸入且是到其像$F(M)$（赋予$N$的子空间拓扑）的同胚，称为光滑嵌入（embedding）。
子集$S\subset N^n$称为$k$维嵌入子流形，若每个$p\in S$有$N$的坐标卡$(U,y)$满足
\[
 y(U\cap S)=y(U)\cap(\R^k\times\{0\}).
\]
这里$S$取子空间拓扑及这些卡诱导的光滑结构。浸入子流形则须给出本身的流形结构与单射浸入$\iota:S\to N$，不要求其拓扑等于像的子空间拓扑。
\end{definition}
\begin{theorem}[局部正规形与嵌入的像]\label{thm:immersion-submersion}
浸入在每一点附近可写为$(u^1,\ldots,u^m)\mapsto(u^1,\ldots,u^m,0,\ldots,0)$；浸没可写为$(u^1,\ldots,u^m)\mapsto(u^1,\ldots,u^n)$。每个浸入限制到充分小的源邻域都是嵌入。光滑嵌入的像为嵌入子流形，反之嵌入子流形的包含映射为光滑嵌入。
\end{theorem}
\begin{proof}
在流形坐标中使用定理\ref{thm:constant-rank}。浸入的秩为源维数，浸没的秩为目标维数，故得到相应正规形。浸入正规形在小开盒上有连续逆（投影到前$m$个坐标），其逆也光滑，因而限制为嵌入。

若$F$为嵌入，在$p$处先选上述源邻域$V$与目标正规形邻域$U$。因$F:M\to F(M)$是同胚，$F(V)$在$F(M)$相对开，故有目标开集$O$使$O\cap F(M)=F(V)$。把$U$缩到$U\cap O$，排除了$V$以外的其他像分支。再缩成适当的坐标开盒，像便恰为其中的坐标平面。这证明整个像满足子流形定义。反过来，子流形坐标之间的过渡是环境光滑过渡在坐标平面上的限制，故光滑；子空间的Hausdorff性和第二可数性从环境遗传。包含映射的矩阵为$\binom{I_k}{0}$，且本来就是到其像的同胚，故为嵌入。
\end{proof}
\begin{definition}[正则值]
对光滑$F:M^m\to N^n$，若$p$处$\dd F_p$满射，称$p$为正则点。若$F^{-1}(q)$中每一点都是正则点，称$q\in N$为正则值（regular value）。空原像时该条件真；此时下述结论按空子流形理解。
\end{definition}
\begin{theorem}[正则值定理与切空间]\label{thm:regular-value}
若$q$为$F:M^m\to N^n$的正则值且$F^{-1}(q)\ne\varnothing$，则$S=F^{-1}(q)$是$m-n$维嵌入子流形。在包含映射的微分所给的识别下，
\[
 T_pS=\ker\dd F_p\subset T_pM\qquad(p\in S).
\]
\end{theorem}
\begin{proof}
满秩小行列式在邻域上仍非零，因此$F$在$p$附近是浸没。应用上一命题并令$q$的坐标为零，$F$的表达为投影，原像就是前$n$个坐标为零的平面。由此得到子流形卡，包含映射的微分将其切空间送到前$n$个分量为零的空间，恰为该投影的核。各卡的光滑兼容性由环境换卡保证，故局部构造拼成所断言的子流形。
\end{proof}
\begin{example}[球面、矩阵群与奇异方程]
$F(x)=|x|^2$在$F^{-1}(R^2)$（$R>0$）上微分为$v\mapsto2\langle x,v\rangle$，故球面维数为$n-1$且$T_xS_R^{n-1}=x^\perp$。这与习题4.6的直接构造相容。
对$G(A)=A^{\mathsf T}A$，把值域取为对称矩阵空间$\operatorname{Sym}_n$。在$A\in O(n)=G^{-1}(I)$处，
\[
 \dd G_A(H)=A^{\mathsf T}H+H^{\mathsf T}A.
\]
给定对称$B$，取$H=AB/2$便得到$\dd G_A(H)=B$，故$O(n)$是$n(n-1)/2$维子流形，$T_AO(n)=\{AK:K^{\mathsf T}=-K\}$。若误把值域取为所有矩阵，满射条件就不成立。
另一方面，$h(x,y)=xy$的零点集是两坐标轴，原点处不能成为一维子流形：任意小邻域去掉原点都有四条局部分支，而一维流形坐标区间只有两条。$g(t)=t^3$的零点集却是零维子流形，尽管$0$不是正则值；正则值条件是充分条件。
\end{example}
\begin{proposition}[紧源与嵌入]\label{prop:compact-embedding}
从紧流形到Hausdorff流形的单射浸入必为嵌入。
\end{proposition}
\begin{proof}
它是到其像的连续双射。紧空间的闭集仍紧，其像在Hausdorff空间中闭，因此逆映射连续。再用浸入假设即得嵌入。
\end{proof}
\begin{example}[单射浸入的逆不一定连续]\label{ex:irrational-immersion}
把$\R/2\pi\mathbb Z$记为$S^1$，取无理数$\alpha$。映射
\[
 \gamma:\R\longrightarrow S^1\times S^1,\qquad
 t\longmapsto([t],[\alpha t])
\]
为单射浸入：同像迫使$t-s=2\pi k$且$\alpha k$为整数，故$k=0$；速度在局部角坐标中是$(1,\alpha)$。它不是嵌入。事实上，把$0,\alpha,\ldots,N\alpha$的小数部分放进$N$个等分区间，抽屉原理给出$1\leq q\leq N$使$q\alpha$到某整数的距离至多$1/N$。当$N\to\infty$可取子列$q_j\to\infty$，否则有限个非零$q$到整数的正距离不能趋零。因此$\gamma(2\pi q_j)\to\gamma(0)$，而$2\pi q_j$不趋于零；到像的逆不连续。无需预先证明整条轨道稠密，就可识别这一拓扑障碍。
\end{example}
本节的正则值和浸入计算见习题5.1--5.2。一般横截原像定理、Sard定理与嵌入存在定理不在本节证明范围内。


\section{向量场、积分曲线与局部流}
向量场$X$是切丛的光滑截面，局部写作$X=a^i(x)\partial_i$；由第4章的截面判别，$X$光滑等价于各$a^i$光滑。对局部函数$f$，$Xf=a^i\partial_i f$。这与梁书2.2.2以“将光滑函数送到光滑函数”表述的定义相容；逆向可用局部坐标函数和第4章的截断恢复系数。
\begin{definition}[积分曲线与完备性]
开区间$I\subset\R$上的光滑曲线$c:I\to M$若满足$\dot c(t)=X_{c(t)}$，称为$X$的积分曲线（integral curve）。取$0\in I$并指定$c(0)=p$，即为初值问题。若每个初值的极大解都定义在整个$\R$，称$X$完备（complete）。零向量处的常值曲线也算积分曲线。
\end{definition}
在坐标中，这成为$\dot x=a(x)$。以下补证本章所需的ODE工具，光滑依赖只针对有限维初值和参数，不预设无限维隐映射定理。
\begin{lemma}[积分不等式]\label{lem:gronwall}
若连续非负函数$u$在$[0,T]$满足$u(t)\leq A+L\int_0^t u(s)\dd s$，其中$A,L\geq0$，则$u(t)\leq Ae^{Lt}$。
\end{lemma}
\begin{proof}
先令$A>0$，置$v(t)=A+L\int_0^t u(s)\dd s$。则$v'\leq Lv$，所以$(e^{-Lt}v)'\leq0$，得到$v\leq Ae^{Lt}$及所需结论。$A=0$时将右边换为任意$\varepsilon>0$，最后令$\varepsilon\downarrow0$。反向时间用$t\mapsto-t$得到相同估计。
\end{proof}
\begin{theorem}[局部解与参数光滑依赖]\label{thm:ode-local}
设$\Omega\subset\R^n$、$\Lambda\subset\R^\ell$开，$a:\Omega\times\Lambda\to\R^n$光滑。在任意$(p_0,\lambda_0)$附近，有$T>0$、初值邻域$V$与参数邻域$W$，使
\[
 \partial_t u(t,p,\lambda)=a(u(t,p,\lambda),\lambda),\qquad
 u(0,p,\lambda)=p
\]
在$(-T,T)\times V\times W$有唯一局部解，且$u$对全部变量光滑。唯一性指任意两解在共同的含零时间区间上一致。
\end{theorem}
\begin{proof}
\textbf{构造解。}在定义域内选紧乘积$\overline B(p_0,2r)\times\overline W$，令$M,L$分别控制$|a|$与$\|D_xa\|$。取$T$使$TM<r$、$TL<1$（若常数为零，相应限制自动满足）。对$p\in B(p_0,r)$，在连续路径空间中考虑
\[
 E_p=\{v\in C([-T,T],\R^n):\|v-p\|_\infty\leq r\},\quad
 (\mathcal T_{p,\lambda}v)(t)=p+\int_0^t a(v(s),\lambda)\dd s.
\]
$E_p$对一致范数完备：一致Cauchy列逐点收敛于欧氏值函数，一致极限连续，闭球约束也保留。$\mathcal T$保持$E_p$并具有压缩常数$TL$。引理\ref{lem:contraction}的迭代与几何级数证明原样适用此完备空间，得到唯一不动点。积分基本定理使该不动点成为$C^1$解，再由方程对时间逐次求导得到时间光滑性。

\textbf{唯一性与连续依赖。}任意两条从同一点出发的解先限制于共同的小坐标球，差的范数满足引理\ref{lem:gronwall}中$A=0$的估计，故相同；沿共同时间区间重复便得唯一性。若$C$控制$D_\lambda a$，不同数据的解满足
\[
 |u(t,p,\lambda)-u(t,q,\mu)|
 \leq (|p-q|+TC|\lambda-\mu|)e^{L|t|}.
\]
估计同时控制整个时间段，并与$|\partial_tu|\leq M$一起给出联合连续性。

\textbf{一阶参数导数。}将$z=(p,\lambda)$合记，候选导数$J=D_z u$必须满足
\[
 J(t)=P+\int_0^t\{A(s)J(s)+B(s)\}\dd s,
\]
其中$P=(I_n\ 0)$，$A(s)=D_xa(u(s,z),\lambda)$，$B(s)=(0\ D_\lambda a(u(s,z),\lambda))$。线性积分算子在同一时间段仍为压缩，故有唯一连续解$J$。其系数对$z$连续，用上一个差估计的同样证明可知$J$联合连续。
取小增量$h$，记$\Delta u=u(\cdot,z+h)-u(\cdot,z)$。上一步给出$\|\Delta u\|_\infty\leq C_1|h|$。$Da$在所选紧集上一致连续，故积分式的Taylor余项一致为$o(|h|)$。减去$J(t)h$的方程，所得余差$R$满足
\[
 |R(t)|\leq o(|h|)+L\left|\int_0^t |R(s)|\dd s\right|.
\]
分别在正、负时间使用引理，得$\|R\|_\infty=o(|h|)$。因此$D_z u=J$，且该导数连续。

\textbf{高阶导数。}为明确归纳步骤，对$z$的多重指标$\beta$形式微分积分方程：最高阶项总是
\[
 U_\beta(t)=P_\beta+\int_0^t
 \{A(s)U_\beta(s)+R_\beta(s)\}\dd s.
\]
$|\beta|\geq2$时$P_\beta=0$；$R_\beta$是$a$的至多$|\beta|$阶导数与$u$的低于$|\beta|$阶导数的有限乘积之和，来自反复使用乘积和链式法则。假设已得$C^{k-1}$，先用这个线性方程构造每个$|\beta|=k$的连续候选。对$|\beta|=k-1$的方程作一阶差商，所有低阶项按归纳假设收敛；最高阶未知仍只以$A$线性出现。把差商减去候选，用$TL<1$或上述积分估计控制余差，便得差商在紧子邻域上一致收敛到候选。故$u$为$C^k$。从$k=2$归纳得到所有$z$导数。最后$\partial_tu=a(u,\lambda)$及其逐次微分给出所有含时间的混合导数连续，得到联合$C^\infty$性。
\end{proof}
非自治方程$\dot x=b(t,x,\lambda)$也可归入此定理。在扩大的空间中令
\[
 (\dot s,\dot x)=(1,b(s,x,\lambda)).
\]
第一分量的解为$s=s_0+t$，代入第二分量便恢复原方程。这使后续线性变分方程和时变方程可使用同一局部工具。
\begin{theorem}[极大局部流]\label{thm:local-flow}
每个光滑向量场$X$存在唯一光滑映射$\Phi:\mathcal D\to M$，其中$\mathcal D\subset\R\times M$开，且
\[
 \mathcal D\cap(\R\times\{p\})=I_p\times\{p\},\qquad
 \Phi(t,p)=\Phi_t(p)
\]
为从$p$出发的极大积分曲线，$I_p$是含零的开区间。满足$\Phi_0(p)=p$；若$s\in I_p$，则
\[
 I_{\Phi_s(p)}=I_p-s,\qquad
 \Phi_t(\Phi_s(p))=\Phi_{t+s}(p).
\]
固定$t$时，$\Phi_t:M_t\to M_{-t}$为微分同胚，$M_t=\{p:t\in I_p\}$，逆为$\Phi_{-t}$。
\end{theorem}
\begin{proof}
坐标中的定理\ref{thm:ode-local}给出局部解；换卡时链式法则将其送成同一向量场方程的解，唯一性使它们相容。对固定$p$，所有含零解区间的并仍为含零开区间，在交上解相同，故拼成极大解。两条时间平移后的解有相同初值，唯一性和极大性给出区间平移关系及复合法则。

还需证明$\mathcal D$开及光滑性。给定$t_0\in I_p$，连接$0$与$t_0$的解段紧。每个段上点附近都有统一短时间的光滑局部解。沿紧时间段取有限分割，使每一步都在一个这样的邻域中；逐步缩小初始点邻域，连续性保证邻近初值也能完成这些有限步。各步解映射的有限复合对初值光滑，最后一步的时间还可在$t_0$附近变化。因此$(t_0,p)$有乘积开邻域包含于$\mathcal D$，且$\Phi$在那里光滑。负$t_0$完全同理。$M_t$由开性得开，复合法则给出$\Phi_{-t}$是$\Phi_t$的两边逆，且二者光滑。
\end{proof}
\begin{example}[局部流不等于全局群]
在$\R$上$X=x^2\partial_x$的解为$\Phi_t(p)=p/(1-tp)$。当$p>0$，极大区间为$(-\infty,1/p)$；当$p<0$，为$(1/p,\infty)$；$p=0$时为整个$\R$。统一的流定义域是$\{(t,p):1-tp>0\}$。代数表达式在分母为负处也有意义，却不能越过爆破时刻拼成含零的解。习题5.4进一步比较定义域。
\end{example}
\begin{proposition}[紧流形上的完备性]\label{prop:compact-complete}
紧且无边界的光滑流形上，每个光滑向量场完备。
\end{proposition}
\begin{proof}
局部解定理给每点一个初值邻域$V_j$和统一存在时间$(-\varepsilon_j,\varepsilon_j)$。紧致性允许取有限个$V_j$覆盖$M$，令$\varepsilon=\min_j\varepsilon_j>0$。若某极大区间有有限右端$b$，取$t\in I_p$满足$b-t<\varepsilon$。从$\Phi_t(p)$出发仍有长度$\varepsilon$的正向解，与旧解借唯一性拼接后越过$b$，矛盾。左端同理。
\end{proof}
\begin{theorem}[流盒定理]\label{thm:flow-box}
若$X_p\ne0$，则$p$附近有坐标$(s,z^2,\ldots,z^n)$使$X=\partial_s$。
\end{theorem}
\begin{proof}
先选以$p$为零点的坐标，经线性变换使$X_p=\partial_1|_p$。取横截片
\[
 \sigma(z)=(0,z^2,\ldots,z^n),\qquad H(s,z)=\Phi_s(\sigma(z)).
\]
在$(0,0)$处，$\partial_sH=X_p$，其余偏导为$\partial_2,\ldots,\partial_n$，故微分可逆。由逆映射定理缩小为微分同胚。流方程给出$\dd H(\partial_s)=X\circ H$，所以逆坐标中$X=\partial_s$。一维时横截片为单点，同一证明适用。
\end{proof}
\begin{figure}[htbp]
\centering\DGDiagram{flow-box}

\caption{原创示意：$X=\partial_x+x\partial_y$的积分曲线为$y=z+x^2/2$。横截线$x=0$上的$z$标记不同曲线，沿曲线的流参数为$s=x$；右图在坐标$(s,z)=(x,y-x^2/2)$中将其拉直。两图坐标单位相同，非等距换坐标会改变曲线形状。}
\label{fig:flow-box}
\end{figure}
若$X_p=0$，任何坐标变换仍把零向量送为零，不能把它变成非零坐标场。流盒只在非零点成立；它描述局部运动，也不保证轨道在整个流形中闭合或嵌入。


\Needspace{10\baselineskip}
\section{李括号与流的交换}
两个方向上的运动次序可能不同。李括号（Lie bracket）刻画这种差别的一阶无穷小结构，而不是对不同点的切向量直接作差。
\begin{definition}[向量场的李括号]
对同一开集上的光滑向量场$X,Y$，定义
\[
 [X,Y]f=X(Yf)-Y(Xf).
\]
符号顺序固定为先写$XY-YX$。这与梁书2.2.2定义10的对易子约定相同。
\end{definition}
\begin{proposition}[坐标式与代数性质]\label{prop:lie-bracket}
李括号仍为光滑向量场。若$X=a^i\partial_i$、$Y=b^i\partial_i$，则
\[
 [X,Y]=(a^j\partial_jb^i-b^j\partial_ja^i)\partial_i.
\]
它是实双线性、反对称的，并满足Jacobi恒等式以及
\[
 [X,[Y,Z]]+[Y,[Z,X]]+[Z,[X,Y]]=0,
\]
\[
 [fX,gY]=fg[X,Y]+fX(g)Y-gY(f)X.
\]
\end{proposition}
\begin{proof}
展开$X(Yf)$与$Y(Xf)$，二阶项分别为$a^jb^i\partial_j\partial_i f$和$b^ja^i\partial_j\partial_i f$，重命名指标并用混合偏导交换后二者抵消，剩下所列一阶算子。其系数光滑且满足点态Leibniz法则，因此为光滑向量场。定义只涉及函数作用，不依赖坐标，局部公式自动相容。反对称性和实双线性直接由算子交换子得到。把Jacobi左边展开为十二个三重复合算子，每个排列各以相反符号出现一次，故为零。最后一式对任意函数$h$展开$fX(gYh)-gY(fXh)$，收集$XYh-YXh$及系数的一阶导数即可。
\end{proof}
李括号通常不是函数双线性的；例如$[\partial_x,x\partial_y]=\partial_y$。特别地，它不是只由$X_p,Y_p$决定的逐点运算：在原点，$Y=x\partial_y$和零场取值相同，却与$\partial_x$有不同括号。习题5.3核查这些区别。
\Needspace{10\baselineskip}
\begin{proposition}[相关向量场的括号]\label{prop:related-bracket}
设$F:M\to N$光滑，$X_1,X_2$与$Y_1,Y_2$分别$F$相关，即$\dd F(X_i)=Y_i\circ F$。则$[X_1,X_2]$与$[Y_1,Y_2]$也$F$相关。特别地，微分同胚保持李括号。
\end{proposition}
\begin{proof}
对$N$上的局部函数$h$，相关性等价于$X_i(h\circ F)=(Y_i h)\circ F$。连续使用两次并相减，得到
\[
 [X_1,X_2](h\circ F)=([Y_1,Y_2]h)\circ F.
\]
局部函数的微分区分切向量，因此得到所需相关性。对微分同胚可写成$F_*[X_1,X_2]=[F_*X_1,F_*X_2]$。
\end{proof}
\begin{lemma}[沿流比较向量场的符号]\label{lem:bracket-flow}
设$\Phi$为$X$的局部流。将$Y_{\Phi_t(p)}$拉回到固定的$T_pM$，记
\[
 A_t(p)=\dd(\Phi_{-t})_{\Phi_t(p)}Y_{\Phi_t(p)}.
\]
则在定义域内
\[
 \frac{\dd}{\dd t}A_t(p)
 =\dd(\Phi_{-t})_{\Phi_t(p)}[X,Y]_{\Phi_t(p)}.
\]
因此在$t=0$导数为$[X,Y]_p$；若改用固定目标点的推前$(\Phi_t)_*Y$，零时导数的符号相反。
\end{lemma}
\begin{proof}
先在$p$处取坐标，记$X=a^i\partial_i,Y=b^i\partial_i$。流方程与光滑性给出
\[
 \Phi_t(p)=p+ta(p)+o(t),\quad
 D\Phi_t(p)=I+tDa(p)+o(t).
\]
后一式由对初值求导的变分方程得到。$D\Phi_{-t}(\Phi_t(p))=D\Phi_t(p)^{-1}$，故$A_t(p)$的展开为
\[
 (I-tDa(p))(b(p)+tDb(p)a(p))+o(t)
 =b(p)+t(Db\,a-Da\,b)(p)+o(t).
\]
这证明零时公式。任意$t$处用流复合律写
\[
 A_{t+h}(p)=\dd(\Phi_{-t})_{\Phi_t(p)}
 \bigl(\dd(\Phi_{-h})_{\Phi_h(q)}Y_{\Phi_h(q)}\bigr),\qquad q=\Phi_t(p),
\]
括号内对$h$在零求导，再套零时公式便得一般式。固定$q$的推前是$\dd\Phi_t|_{\Phi_{-t}(q)}Y_{\Phi_{-t}(q)}=A_{-t}(q)$，因此符号反转。
\end{proof}
\Needspace{12\baselineskip}
\begin{theorem}[括号为零与局部流交换]\label{thm:commuting-flows}
对开集上的光滑场$X,Y$，下列条件等价：$[X,Y]=0$；$X$的局部流保持$Y$；两者的局部流在每点附近、充分小的时间内交换，即
\[
 \Phi_s\circ\Psi_t=\Psi_t\circ\Phi_s.
\]
这里每个复合须有定义；结论不声称在未经检查的任意大时间矩形上成立。
\end{theorem}
\begin{proof}
若$[X,Y]=0$，引理\ref{lem:bracket-flow}说明$A_s(p)$为常量$Y_p$，即$\dd\Phi_s(Y_p)=Y_{\Phi_s(p)}$。这是流保持$Y$。
若保持，固定足够小的$s$，曲线$t\mapsto\Phi_s(\Psi_t(p))$的导数等于其所在点的$Y$值，初值为$\Phi_s(p)$。ODE唯一性使它等于$\Psi_t(\Phi_s(p))$，故局部交换。反之，若交换，对$t$在零求导即得保持性，再由引理在$s=0$求导得到$[X,Y]=0$。所有推导可先缩小到统一的小时间和空间邻域，避免定义域越界。
\end{proof}
\begin{example}[交换矩形的剩余位移]\label{ex:bracket-rectangle}
在$\R^3$取$X=\partial_x$、$Y=\partial_y+x\partial_z$。其流为
\[
 \Phi_s(x,y,z)=(x+s,y,z),\qquad
 \Psi_t(x,y,z)=(x,y+t,z+xt).
\]
依次沿$X$走$s$、沿$Y$走$t$、沿$X$走$-s$、沿$Y$走$-t$，终点为$(x,y,z+st)$。故与约定$[X,Y]=\partial_z$一致。次序倒置会改变符号；只写“四边绕行”而不写次序不足以确定括号方向。
\end{example}
\begin{figure}[htbp]
\centering\DGDiagram{bracket-loop}

\caption{原创示意：例\ref{ex:bracket-rectangle}从原点取$s=t=1$，按箭头依次走$X,Y,-X,-Y$，终点$q=(0,0,1)$。灰虚线表示剩余位移$st\partial_z$；空间采用斜投影，画面中的线段长度不直接代表欧氏距离。}
\label{fig:bracket-loop}
\end{figure}


\section{Frobenius定理及其证明}
曲面在每一点给出一个切平面。反过来，先指定随点变化的切平面，是否总能找到与它们处处相切的曲面？一维的答案由流盒给出，高维则需要相容条件。
\begin{definition}[光滑分布与积分子流形]
秩$k$的光滑分布（smooth distribution）$D$是在每点指定$k$维子空间$D_p\subset T_pM$，并要求每点附近有光滑场$X_1,\ldots,X_k$逐点构成$D$的基。写$\Gamma(D)$表示取值于$D$的光滑向量场，局部版本同样使用。若对所有局部截面$X,Y$有$[X,Y]\in\Gamma(D)$，称$D$对合（involutive）。
维数为$k$的浸入$\iota:L\to M$若满足$\dd\iota_q(T_qL)=D_{\iota(q)}$，称为积分流形；若$\iota$还单射，称为积分子流形。本节的局部积分片均是嵌入子流形。
\end{definition}
局部标架存在时，任意截面都可唯一写作$f^iX_i$，系数光滑：选标架分量矩阵的一个非零$k$阶小行列式，在邻域上求逆即可。由命题\ref{prop:lie-bracket}，检验对合性只需检查$[X_i,X_j]$是否落在$D$，无需逐一检查所有截面。
\begin{theorem}[Frobenius局部坐标定理]\label{thm:frobenius}
设$D$是$M^n$上常秩$k$的光滑分布，$0\leq k\leq n$。下列三条件等价：
\begin{enumerate}[label=(\roman*),nosep]
\item $D$对合；
\item 每点有局部坐标$(u^1,\ldots,u^k,v^1,\ldots,v^{n-k})$，使$D$由$\partial_{u^1},\ldots,\partial_{u^k}$张成；
\item 每点都有经过它的$k$维积分流形。
\end{enumerate}
在(ii)的乘积坐标盒内，固定$v$得到嵌入积分片；任意积分流形在每个点的充分小邻域上落在唯一这样的片中。
\end{theorem}
\begin{proof}
\textbf{(i)推出(ii)：先改造标架，再造坐标。}取$p$附近标架$Z_1,\ldots,Z_k$及环境坐标$x^1,\ldots,x^n$。重排$x$可使标架的前$k$行分量矩阵在$p$可逆，并缩小邻域使其处处可逆。用该矩阵的逆对标架作光滑线性组合，得到
\[
 Y_i=\partial_{x^i}+\sum_{\alpha=k+1}^n
 a_i^\alpha(x)\partial_{x^\alpha},\qquad 1\leq i\leq k.
\]
于是$D$在每点由这些$Y_i$张成。因$Y_i(x^j)=\delta_i^j$（$j\leq k$），有$[Y_i,Y_j](x^\ell)=0$（$\ell\leq k$），即该括号的前$k$个分量为零。另一方面，对合性使$[Y_i,Y_j]=c^rY_r$。比较前$k$个分量得每个$c^r=0$，故$Y_i$两两交换。此处原标架$Z_i$未必交换，不能直接把它们当成坐标基。

平移坐标使$p=0$，取横截片$\sigma(v)=(0,\ldots,0,v^1,\ldots,v^{n-k})$。设$\Phi^{(i)}$为$Y_i$的局部流，定义
\[
 H(u,v)=\Phi^{(1)}_{u^1}\circ\cdots\circ
 \Phi^{(k)}_{u^k}(\sigma(v)).
\]
由定理\ref{thm:local-flow}可缩小到统一的参数盒使复合光滑有定义。定理\ref{thm:commuting-flows}使各流局部交换且相互保持相应向量场，从而
\[
 \dd H(\partial_{u^i})=Y_i\circ H.
\]
在$(0,0)$处，$DH$的列为$Y_1(p),\ldots,Y_k(p),\partial_{x^{k+1}},\ldots,\partial_{x^n}$；其分量矩阵具有对角块$I_k,I_{n-k}$，所以可逆。逆映射定理给出局部坐标$H^{-1}$，在其中$D$恰由前$k$个坐标场张成。$k=0$时任何坐标均可，$k=n$时横截片是单点，构造仍成立。

\textbf{(ii)推出(i)与(iii)。}坐标场两两交换，$D$的任意两个截面由前$k$个坐标场线性组合，括号的乘积公式说明结果仍为它们的组合，因此对合。固定$v=c$的切空间由这些坐标场张成，所以是$k$维嵌入积分片，每点都在一个片上。

\textbf{(iii)推出(i)。}给定$p$及经过它的积分浸入$\iota:L\to M$，在其某个原像点附近限制为嵌入，以子流形卡表示像为$(w^1,\ldots,w^k,0,\ldots,0)$。$D$的任意局部截面$X,Y$在这个像上都相切，所以其法向分量在像上为零。法向分量沿切方向的导数也为零；将此代入李括号坐标式，$[X,Y]$的全部法向分量仍为零。在$p$处便得$[X,Y]_p\in D_p$。因$p$任意，$D$对合。

\textbf{局部唯一性。}在(ii)的坐标盒中，积分浸入$\iota$满足$\dd(v^\alpha\circ\iota)=0$。限制到源上连通的凸坐标邻域，沿线段积分可知各$v^\alpha\circ\iota$为常数。故像落在唯一的$v=c$积分片中。再比较维数，$u\circ\iota$的微分为同构，由逆映射定理该局部像是积分片的开子集。因此经过同一点的积分片在缩小后相同，即局部积分片的芽唯一。
\end{proof}
\begin{corollary}[局部第一积分]\label{cor:first-integrals}
对合的常秩$k$分布在每点附近可写为
\[
 D=\ker\dd v,\qquad v=(v^1,\ldots,v^{n-k}),\quad
 \operatorname{rank}\dd v=n-k.
\]
反之任意浸没$v$的核分布都是对合的，其纤维的连通局部片是积分片。
\end{corollary}
\begin{proof}
Frobenius坐标中投影到$v$即可。反向在浸没坐标中核由其余坐标场张成；或者对$X,Y\in\ker\dd v$计算$[X,Y]v^\alpha=X(Yv^\alpha)-Y(Xv^\alpha)=0$。核的光滑标架由投影正规形给出，积分片由正则值定理得到。
\end{proof}
\begin{example}[可积与不可积的平面场]
在$\R^3$上，$D=\operatorname{span}(\partial_x+y\partial_z,\partial_y+x\partial_z)$是$\ker\dd(z-xy)$，两基场的括号为零。其积分曲面是$z-xy=c$；局部可使用$(u,v,w)=(x,y,z-xy)$作Frobenius坐标。
与此相对，例\ref{ex:bracket-rectangle}中的
\[
 E=\operatorname{span}(\partial_x,\partial_y+x\partial_z)
\]
为光滑常秩二分布，然而$[X,Y]=\partial_z\notin E$：若$\partial_z=a\partial_x+b(\partial_y+x\partial_z)$，比较$x,y$分量迫使$a=b=0$，与$z$分量矛盾。故任意点都不存在二维积分曲面。四段流的剩余位移显示了无法封闭成曲面的障碍。
\end{example}
\begin{figure}[htbp]
\centering\DGDiagram{integral-plaques}

\caption{原创示意：$z-xy=c$的三个局部积分片（$c=-1.8,0,1.8$，$|x|,|y|\leq0.75$）。每片的切空间由$\partial_x+y\partial_z$与$\partial_y+x\partial_z$张成；右图为$(u,v,w)=(x,y,z-xy)$坐标中的平行片$w=c$。采用斜投影，仅表达局部叶片与横截参数，不表示等距变换。}
\label{fig:integral-plaques}
\end{figure}
本定理要求常秩与局部光滑标架。若只指定任意点态子空间，或允许秩跳跃，就不能直接套用。局部积分片可作为叶状结构的局部模型；最大连通叶的整体构造、叶空间的拓扑与奇异分布理论另属后续专题，本章不据局部唯一性宣称叶总为嵌入或叶空间总为流形。例\ref{ex:irrational-immersion}已经提示了全局拓扑的差别。


\Needspace{15\baselineskip}
\section{子流形与局部流习题及详细解答}
习题5.3、5.4以梁书上册第二版2.2.2的文内命题和选读说明为出发点，增加坐标计算与定义域核查；并非原书章末题号。其余六题为原创补充。来源位置与双页码见\texttt{references/EXERCISES.md}，全部解答由本工程编写。

\Needspace{10\baselineskip}
\subsection*{习题5.1：双约束与切空间}
设$F:\R^3\to\R^2$，$F(x,y,z)=(x^2+y^2+z^2,z)$。判断$(1,c)$何时为非空纤维的正则值，并求纤维及其切空间。解释$c=\pm1$和$|c|>1$。
\begin{proof}[解答]
纤维满足$z=c$、$x^2+y^2=1-c^2$。当$|c|<1$时为水平圆，$DF$的两行是$(2x,2y,2c)$与$(0,0,1)$，因$(x,y)\ne(0,0)$而线性无关，故为正则值。切空间为
\[
 \{(a,b,d):d=0,\ xa+yb=0\}
 =\operatorname{span}\{(-y,x,0)\}.
\]
当$c=\pm1$时纤维是单点$(0,0,c)$，其本身是零维子流形，但$DF$秩为一，值不是正则值，不能用$m-n=1$推断其维数。$|c|>1$时纤维为空，按定义仍是正则值。可见“正则值”不保证纤维非空，奇异值也不强迫纤维不是子流形。
\end{proof}

\Needspace{9\baselineskip}
\subsection*{习题5.2：自交浸入与紧源}
证明$\gamma:S^1\to\R^2$、$\gamma([t])=(\sin t,\sin2t)$为浸入而非嵌入。说明为什么这不违背紧源命题；再证明其限制到充分小的参数弧上为嵌入。
\begin{proof}[解答]
导数为$(\cos t,2\cos2t)$；若$\cos t=0$，则$\cos2t=-1$，所以两分量不会同时为零，故为浸入。$[0]\ne[\pi]$却同映到$(0,0)$，所以不单射，也就不是嵌入。命题\ref{prop:compact-embedding}同时要求单射，不能只用紧致性。对任意$t_0$，至少一个分量的导数非零，该分量在小参数弧上由逆映射定理给出局部参数，因而$\gamma$在该弧上单射且到像的逆连续、光滑。两个分支在环境中相遇，不妨碍源上每个分支独立地局部嵌入。
\end{proof}

\Needspace{10\baselineskip}
\subsection*{习题5.3：对易子、坐标基与函数系数}
依据梁书印刷页28定义10及定理2-2-7，证明坐标场$\partial_x,\partial_y$的括号为零。在$\R^2$取$X=\partial_x,Y=x\partial_y$，计算$[X,Y]$，并在原点比较$Y$与零场。向量场的李括号能否作为$T_0\R^2$上的双线性运算？
\begin{proof}[解答]
对局部光滑$f$，$[\partial_x,\partial_y]f=\partial_x\partial_yf-\partial_y\partial_xf=0$。一般式则给$[\partial_x,x\partial_y]=\partial_x(x)\partial_y=\partial_y$。原点处$Y_0=0$，但$[X,Y]_0=\partial_y|_0$，而$[X,0]_0=0$。所以括号的值还读取向量场邻域内的变化，不能只由两个点向量定义；若强行定义为点态双线性运算，与零向量的括号应为零，立即矛盾。这里的双线性是对实常数的双线性，不是对光滑函数的双线性。
\end{proof}

\Needspace{10\baselineskip}
\subsection*{习题5.4：最大积分曲线的时间区间}
参照梁书印刷页28--29定理2-2-8及选读2-2-3的局部唯一性与不可延拓说明，求$X=x^2\partial_x$各初值的极大积分曲线。限制到$M=(-1,1)$时，从$p=0$与$p=1/2$出发的极大时间区间分别是什么？
\begin{proof}[解答]
$p=0$时由唯一性得到恒零解，区间为$\R$。$p\ne0$时在解不为零的区间积分$(1/x)'=-1$，得$x(t)=p/(1-pt)$。正$p$的含零区间为$(-\infty,1/p)$，负$p$为$(1/p,\infty)$；有限端点处$|x|\to\infty$，不能连续延拓为实线值解。到无穷远的另一端没有有限时间端点，故这些区间极大。
在$M=(-1,1)$上，$p=0$仍为恒零全局解。$p=1/2$时$x(t)=1/(2-t)$，要求$-1<x(t)<1$并留在含零的连续分支，得到$t<1$，故区间为$(-\infty,1)$。当$t\uparrow1$时趋向被删去的边界点$1$，虽未在$\R$中爆破，却不能继续作为$M$值解。这说明极大区间同时依赖方程和所在流形。
\end{proof}

\Needspace{10\baselineskip}
\subsection*{习题5.5：显式流盒}
对$X=\partial_x+x\partial_y$求流，构造全平面上的流盒坐标，并证明这些坐标确实使$X$变为一个坐标场。是否由此可知该换坐标保持长度？
\begin{proof}[解答]
从$(x_0,y_0)$出发，$x=x_0+t$，积分$\dot y=x_0+t$得
\[
 \Phi_t(x_0,y_0)=(x_0+t,y_0+x_0t+t^2/2).
\]
令$s=x,z=y-x^2/2$，其逆为$(x,y)=(s,z+s^2/2)$，两者全局光滑，Jacobi行列式恒为一。$X(s)=1$、$X(z)=x-x=0$，所以新坐标中$X=\partial_s$。然而在$(s,z)$的坐标下，欧氏度量为$\dd s^2+(\dd z+s\dd s)^2$，并非$\dd s^2+\dd z^2$，所以换坐标不保持以标准分量计算的长度。流盒结论属于光滑结构，未给出等距性。
\end{proof}

\Needspace{10\baselineskip}
\subsection*{习题5.6：可积平面场与第一积分}
对$D=\operatorname{span}(\partial_x+y\partial_z,\partial_y+x\partial_z)$，求非恒定第一积分和经过$(a,b,c)$的局部积分片；验证两个生成场的流交换。
\begin{proof}[解答]
置$w=z-xy$，则$\dd w=(-y,-x,1)$处处非零，且两生成场均使$w$的导数为零。两场逐点独立、核为二维，故$D=\ker\dd w$。经过指定点的片为$z-xy=c-ab$，切空间等于该核。
第一生成场的流为$\Phi_s(x,y,z)=(x+s,y,z+sy)$，第二场的流为$\Psi_t(x,y,z)=(x,y+t,z+tx)$。两个复合均为
\[
 (x+s,y+t,z+sy+tx+st),
\]
故全局交换。第一积分把局部曲面方程与流的相容性联系起来，但不因此将所有分布都视为某个全局函数的核。
\end{proof}

\Needspace{10\baselineskip}
\subsection*{习题5.7：不可积的混合偏导障碍}
对$E=\operatorname{span}(\partial_x,\partial_y+x\partial_z)$，除李括号判别外，直接证明不存在二维积分曲面。说明为何局部曲面可以先假设为图像$z=h(x,y)$。
\begin{proof}[解答]
若有积分浸入，其切平面投影到$(x,y)$平面是线性同构：两生成向量的投影是标准基。因此对积分浸入复合该投影应用逆映射定理，局部就能把曲面写成图像。图像切空间为$\operatorname{span}(\partial_x+h_x\partial_z,\partial_y+h_y\partial_z)$，与$E$相同且$x,y$投影固定时，唯一提升给出$h_x=0$、$h_y=x$。于是$h_{xy}=0$而$h_{yx}=1$，违反光滑函数混合偏导交换，矛盾。这个证明同时解释了为什么不能凭图像直觉省略局部投影可逆这一步。
\end{proof}

\Needspace{10\baselineskip}
\subsection*{习题5.8：对合标架未必是坐标标架}
在$\R^2$取$Z_1=\partial_x,Z_2=e^x\partial_y$。证明其张成的分布对合，但不存在局部坐标使这两个指定向量场同时成为两个坐标基向量。给出一个光滑标架变换，将它们改成坐标标架。
\begin{proof}[解答]
两场处处独立，张成$T\R^2$，该分布当然对合；也可计算$[Z_1,Z_2]=e^x\partial_y=Z_2$并用标架判别。若两指定场同时为坐标基，则它们的括号应为零，与该计算矛盾。改取$Y_1=Z_1,Y_2=e^{-x}Z_2=\partial_y$，得到原坐标的基，二者交换。Frobenius只保证分布有一个坐标标架，并不承诺输入的任意标架都可直接成为坐标基；正文中的归一化步骤正是为解决这个差别。
\end{proof}

